\hnheader[Warum K-Means konvergiert -- und Code für K-Means, hierarchisches Clustering und Gauß-Gemische.][Jeder Schritt senkt $J$: zuweisen minimiert, mitteln minimiert]{K-Means \& Clustering:}{Vertiefung}{Herleitung und Code}
\hnsrc{main\_notes.pdf Kap.\ 10 (K-Means), notes/cs229-notes7a.pdf; hand-notes/ML Notes.pdf (K-Means, Hierarchical Clustering); Herleitung ausformuliert; Code: code/sk\_cluster.py (real ausgeführt)}

\begin{hngrid}
%
\begin{hnp}[raster multicolumn=3,acc=hnorange]{1}{K-Means: $J$ sinkt in jedem Schritt}
\hnleg{$J$ Kosten (Summe der quadrierten Abstände zum Zentroid), $c^{(i)}$ Cluster von Punkt $i$, $\mu_k$ Zentroid von Cluster $k$, $n_k$ Anzahl Punkte in Cluster $k$.}
\hnwords{Zwei Schritte, die $J$ nie erhöhen: jeden Punkt dem nächsten Zentrum zuweisen, dann jedes Zentrum auf den Mittelwert seiner Punkte setzen.}
\textbf{Zuweisung}: für festes $\mu$ minimiert $c^{(i)}=\arg\min_k\|x^{(i)}-\mu_k\|^2$ jeden Summanden $\Rightarrow J\downarrow$.\\
\textbf{Update}: für festes $c$ gilt
\begin{align*}
\nabla_{\mu_k}\!\!\sum_{i:c^{(i)}=k}\!\|x^{(i)}-\mu_k\|^2&=-2\!\!\sum_{i:c^{(i)}=k}\!(x^{(i)}-\mu_k)=0\\
\Rightarrow\;\mu_k&=\tfrac1{n_k}\textstyle\sum_{i:c^{(i)}=k}x^{(i)}&&\why{Hesse $2n_kI\succ0$}
\end{align*}
$J\ge0$ nicht steigend, endlich viele Zuweisungen $\Rightarrow$ Konvergenz (\emph{lokales} Minimum, nichtkonvex).
\end{hnp}
%
\begin{hnex}[raster multicolumn=3]{Beispiel: Der Mittelwert minimiert}
Punkte $\{1,2,9\}$, ein Cluster, Zentroid $\mu$.\\
\hnsol\ $J(\mu)=\sum(x-\mu)^2$: $J(2)=1+0+49=50$, $J(\mu{=}4)=9+4+25=\ora{38}$ (Mittel), $J(6)=25+16+9=50$.\\
Das Minimum liegt beim arithmetischen Mittel $4$ (numerisch geprüft).
\end{hnex}
%
\hnsnip[hnpurple]{sk_cluster}{K-Means (Inertia, Silhouette), hierarchisches Clustering, Gauß-Gemisch}
%
\begin{hnp}[raster multicolumn=6,acc=hnpurple,colback=hnlilac!30]{?}{Selbsttest: kannst du das herleiten?}
\hnquiz{Warum sinkt $J$ im Zuweisungsschritt?}{Jedes $x^{(i)}$ wählt den nächsten Zentroid, das minimiert seinen Summanden.}{Warum ist der Zentroid der Mittelwert?}{Ableiten und Null setzen: $\sum(x-\mu)=0$, Hesse $2n_kI\succ0$.}{Warum nur ein lokales Minimum?}{$J$ ist nichtkonvex; das Ergebnis hängt von der Initialisierung ab.}
\end{hnp}
%
\begin{hnbanner}[raster multicolumn=6]
„Zuweisen, mitteln, wiederholen -- und jedes Mal wird es besser.“
\end{hnbanner}
\end{hngrid}
